% =====================================================================
%  Capitolo 2 — Premesse: stato dell'arte e lavoro del 2022
% =====================================================================
\chapter{Premesse}
\label{cap:premesse}

Questo capitolo colloca il lavoro nel contesto della letteratura e ne fissa il punto di
partenza. Le prime sezioni ripercorrono le quattro linee di ricerca su cui la tesi si appoggia:
segmentazione semantica, attenzione efficiente, rappresentazioni discrete e apprendimento
auto-supervisionato. Le ultime riassumono la tesi del 2022 e ne propongono una lettura critica,
da cui discendono le scelte di metodo dei capitoli successivi.

\section{Segmentazione semantica con reti profonde}
\label{sec:stato-segmentazione}

Le \emph{Fully Convolutional Networks} (FCN) di Long et al.~\cite{long2015fcn} hanno mostrato
che una rete di classificazione può produrre una predizione densa sostituendo gli strati
completamente connessi con convoluzioni e ricampionando le mappe di attivazione alla risoluzione
dell'ingresso. Il limite principale di questo schema è la perdita di dettaglio spaziale dovuta
al sottocampionamento progressivo. Due famiglie di soluzioni lo affrontano in modo diverso.

\begin{itemize}
  \item Le architetture \textbf{encoder-decoder} con connessioni di salto, come
  U-Net~\cite{ronneberger2015unet}, recuperano il dettaglio combinando, a ogni scala del
  decoder, le mappe ad alta risoluzione dell'encoder. Con un encoder residuale~\cite{he2016resnet}
  si ottengono le varianti \emph{ResU-Net}, che possono riusare pesi pre-addestrati su
  ImageNet~\cite{deng2009imagenet}.
  \item Le architetture con \textbf{convoluzioni dilatate} (\emph{atrous}), come
  DeepLab~\cite{chen2018deeplab}, allargano il campo recettivo senza ridurre la risoluzione,
  a un costo di memoria maggiore.
\end{itemize}

Più di recente gli encoder di tipo transformer, come ViT~\cite{dosovitskiy2021vit} e Swin
Transformer~\cite{liu2021swin}, e i decoder leggeri come quello di
SegFormer~\cite{xie2021segformer} hanno spostato lo stato dell'arte, al prezzo di un costo
computazionale e di un fabbisogno di dati maggiori.

\section{Segmentazione di immagini aeree e il benchmark ISPRS}
\label{sec:stato-aereo}

Il benchmark ISPRS~\cite{rottensteiner2012isprs} ha reso confrontabili i metodi di
classificazione di scene urbane aeree. Tra i primi lavori con reti profonde,
Sherrah~\cite{sherrah2016fcn} e Marmanis et al.~\cite{marmanis2016ensemble} hanno applicato FCN
e loro ensemble. Audebert et al.~\cite{audebert2018beyond} hanno mostrato il vantaggio di
fondere l'informazione altimetrica (DSM/nDSM) con quella spettrale, sia con una fusione precoce
(come canale aggiuntivo) sia con una fusione tardiva tra due rami. Lo stesso gruppo di autori di
\maresunet{} ha poi proposto UNetFormer~\cite{wang2022unetformer}, un decoder con attenzione
globale-locale per scene urbane. MMSegmentation~\cite{mmseg2020} ha consolidato uno split di
riferimento per Potsdam e Vaihingen, adottato da gran parte della letteratura recente, che questa
tesi segue (\cref{sec:split}).

Due caratteristiche distinguono il dominio dalle immagini naturali e guidano le scelte di
questo lavoro. La prima è l'\textbf{assenza di un'orientazione canonica}: una scena aerea
ruotata di \ang{90} o riflessa resta una scena plausibile, quindi le trasformazioni del gruppo
diedrale $D_4$ sono augmentation \enquote{gratuite}, che non alterano le etichette. La seconda è
la \textbf{disponibilità di modalità ausiliarie} co-registrate, come il DSM e l'nDSM, che
portano un'informazione (l'altezza) non ricavabile dal solo colore.

\section{Attenzione efficiente}
\label{sec:stato-attenzione}

L'attenzione \emph{dot-product}~\cite{vaswani2017attention} pesa ogni posizione in base alla
somiglianza tra la sua \emph{query} e le \emph{key} di tutte le altre. Nella visione artificiale è stata
introdotta come blocco \emph{non-local}~\cite{wang2018nonlocal} e, per la segmentazione, da
DANet~\cite{fu2019danet}, che affianca un'attenzione sulle posizioni (PAM) a una sui canali
(CAM). Il costo di memoria e di calcolo cresce però con il quadrato del numero di posizioni
$N=HW$, e questo ne limita l'uso alle mappe a bassa risoluzione.

Le \emph{attenzioni lineari} riscrivono il prodotto $\softmax(QK\T)V$ come $\phi(Q)\bigl(\phi(K)\T
V\bigr)$ per un'opportuna mappa di caratteristiche $\phi$, e sfruttano l'associatività del
prodotto matriciale per ridurre il costo a $\bigO{N}$~\cite{katharopoulos2020linear}. Il
\emph{Linear Attention Mechanism} (LAM) di Li et al.~\cite{li2020lam} ottiene lo stesso risultato
con uno sviluppo di Taylor al primo ordine dell'esponenziale, applicato a query e key
normalizzate. \maresunet{}~\cite{li2022maresunet} inserisce un blocco PAM lineare + CAM in ogni
connessione di salto di una ResU-Net. La derivazione completa è nella
\cref{sec:attenzione-lineare}.

\section{Rappresentazioni discrete}
\label{sec:stato-discrete}

Gli autoencoder variazionali~\cite{kingma2014vae} apprendono una rappresentazione latente
continua regolarizzata verso una distribuzione a priori. Il VQ-VAE~\cite{vandenoord2017vqvae}
sostituisce lo spazio latente continuo con un insieme finito di vettori (\emph{codebook}): ogni
vettore prodotto dall'encoder viene sostituito dal codice più vicino. L'operazione non è
derivabile e il gradiente viene propagato con lo stimatore
\emph{straight-through}~\cite{bengio2013straight}. Le rappresentazioni discrete sono state alla
base di modelli generativi ad alta fedeltà, come VQ-VAE-2~\cite{razavi2019vqvae2} e
VQGAN~\cite{esser2021taming}. Il loro problema pratico più noto è il \emph{collasso del
codebook}: pochi codici vengono usati e gli altri restano \enquote{morti}. Le contromisure più
diffuse sono l'aggiornamento per medie mobili esponenziali (EMA), l'inizializzazione dai dati e
il riavvio dei codici inutilizzati~\cite{lancucki2020robust}. Per la segmentazione, il codebook
agisce come un \emph{information bottleneck}~\cite{tishby2000ib}: limita la quantità di
informazione che attraversa il collo di bottiglia, con un possibile effetto regolarizzante.
L'idea richiama i \emph{visual words} di OBoW~\cite{gidaris2021obow}, in cui le feature vengono
assegnate alla parola più vicina di un vocabolario.

\section{Apprendimento auto-supervisionato}
\label{sec:stato-ssl}

L'apprendimento auto-supervisionato (\emph{self-supervised learning}, SSL) costruisce un compito
ausiliario a partire dai soli dati, senza etichette, e usa l'encoder così addestrato come punto
di partenza per il compito a valle. La \cref{tab:ssl-famiglie} riassume le famiglie principali
nel contesto di questa tesi. La sfida comune a tutti i metodi che confrontano due viste della
stessa immagine è il \emph{collasso}: la soluzione banale in cui la rete produce lo stesso
output per ogni ingresso.

\begin{table}[t]
  \centering
  \small
  \caption[Famiglie di metodi auto-supervisionati]{Famiglie di metodi di apprendimento
    auto-supervisionato per immagini e meccanismo con cui ciascuna evita il collasso.}
  \label{tab:ssl-famiglie}
  \begin{tabular}{|l|p{3.1cm}|p{3.7cm}|l|}
    \hline
    \textbf{Famiglia} & \textbf{Compito} & \textbf{Contro il collasso} & \textbf{Esempi} \\
    \hline
    Contrastiva & distinguere viste della stessa immagine da quelle di altre &
      esempi negativi (batch o coda) & \cite{chen2020simclr,he2020moco} \\
    \hline
    Clustering & predire l'assegnazione a prototipi di un'altra vista &
      vincolo di equipartizione & \cite{caron2020swav} \\
    \hline
    Non contrastiva & predire la rappresentazione di un'altra vista &
      asimmetria: predittore e \emph{stop-gradient}, teacher EMA & \cite{grill2020byol,chen2021simsiam} \\
    \hline
    Auto-distillazione & lo student imita la distribuzione del teacher &
      \emph{centering} e \emph{sharpening} dell'output del teacher, teacher EMA &
      \cite{caron2021dino,li2022esvit} \\
    \hline
    \emph{Bag of visual words} & ricostruire l'istogramma di parole visive del teacher &
      vocabolario aggiornato online, teacher EMA & \cite{gidaris2021obow} \\
    \hline
    Ricostruzione mascherata & ricostruire le porzioni mascherate &
      compito generativo, nessuna coppia di viste & \cite{he2022mae,cong2022satmae} \\
    \hline
  \end{tabular}
\end{table}

L'auto-distillazione discende dalla distillazione della conoscenza di Hinton et
al.~\cite{hinton2015distilling,gou2021kdsurvey}. Il teacher non è più una rete pre-addestrata
più grande, ma una media mobile esponenziale dei pesi dello student, come nel \emph{mean
teacher}~\cite{tarvainen2017meanteacher} e nella media di Polyak~\cite{polyak1992averaging}.
DINO~\cite{caron2021dino} combina questa idea con un'augmentation \emph{multi-crop} (due viste
globali per il teacher, viste globali e locali per lo student) e con la coppia
\emph{centering}/\emph{sharpening}. È stato proposto per i Vision Transformer, ma funziona anche
con encoder ResNet. Le versioni successive, come DINOv2~\cite{oquab2024dinov2}, ne hanno
confermato la scalabilità a grandi collezioni di immagini.

Nel telerilevamento l'SSL è motivato proprio dalla scarsità di etichette~\cite{wang2022ssl_rs}.
Metodi specifici sfruttano la dimensione temporale (SeCo~\cite{manas2021seco}), la
geolocalizzazione (GASSL~\cite{ayush2021gassl}) o la ricostruzione mascherata di immagini
multispettrali (SatMAE~\cite{cong2022satmae}). MARE~\cite{marsocci2021mare} ha applicato OBoW
alla \maresunet{} su Vaihingen e ha osservato che le augmentation cromatiche forti, utili su
ImageNet, non aiutano necessariamente sulle immagini aeree. È l'osservazione all'origine della
\RQ{4}.

\section{La tesi del 2022}
\label{sec:tesi-2022}

La tesi del 2022~\cite{campana2022tesi} ha studiato una griglia di $2\times2\times2$
configurazioni: architettura (\maresunet{}, \maresunetvq{}), encoder (ResNet-18, ResNet-34) e
pre-addestramento (ImageNet supervisionato, DINO su Potsdam). Le scelte sperimentali principali
erano le seguenti.

\begin{itemize}
  \item \textbf{Dati}: Potsdam, ritagli di $256\times256$ non sovrapposti. Erano state rimosse
  dieci aree giudicate troppo omogenee (\tile{2\_10}, \tile{2\_11}, \tile{2\_12}, \tile{2\_13},
  \tile{2\_14}, \tile{3\_10}, \tile{6\_7}, \tile{6\_13}, \tile{7\_13}, \tile{7\_9}); delle 28
  rimanenti, 22 erano usate per il training e 6 per la validation (circa \num{11000} ritagli di
  training).
  \item \textbf{Addestramento}: Adam, learning rate $10^{-3}$, StepLR (passo 15, $\gamma=0{,}4$),
  100 epoche, batch 32, nessuna augmentation in fine-tuning.
  \item \textbf{Pre-addestramento}: DINO con scheduler coseno, \emph{weight decay} da 0,04 a 0,4,
  $\tau_s=0{,}1$ e $\tau_t$ da 0,04 a 0,07.
\end{itemize}

\begin{table}[t]
  \centering
  \small
  \caption[Risultati di test della tesi del 2022]{Risultati di test riportati nella tesi del
    2022 su Potsdam (tab.~4.5 dell'originale), con il protocollo di allora. In grassetto il
    valore migliore di ogni colonna.}
  \label{tab:risultati-2022}
  \begin{tabular}{|l|l|l|c|c|c|}
    \hline
    \textbf{Modello} & \textbf{Encoder} & \textbf{Pre-addestr.} & \textbf{OA} & \textbf{mF1} & \textbf{mIoU} \\
    \hline
    \maresunet{}   & ResNet-18 & ImageNet & 0,791 & \textbf{0,722} & \textbf{0,604} \\
    \maresunet{}   & ResNet-34 & ImageNet & \textbf{0,791} & 0,722 & 0,604 \\
    \maresunet{}   & ResNet-18 & DINO     & 0,741 & 0,657 & 0,530 \\
    \maresunet{}   & ResNet-34 & DINO     & 0,739 & 0,653 & 0,525 \\
    \maresunetvq{} & ResNet-18 & ImageNet & 0,750 & 0,668 & 0,543 \\
    \maresunetvq{} & ResNet-34 & ImageNet & 0,744 & 0,660 & 0,534 \\
    \maresunetvq{} & ResNet-18 & DINO     & 0,702 & 0,593 & 0,465 \\
    \maresunetvq{} & ResNet-34 & DINO     & 0,703 & 0,595 & 0,467 \\
    \hline
  \end{tabular}
\end{table}

I risultati (\cref{tab:risultati-2022}) indicavano che la quantizzazione e il pre-addestramento
DINO peggioravano tutte le metriche e che i loro effetti si sommavano. ResNet-18 e ResNet-34
risultavano equivalenti. Il clutter era la classe peggiore (IoU $\approx 0{,}15$) e alberi e
vegetazione bassa si confondevano. Le metriche di validation (mIoU $\approx 0{,}96$) erano
molto più alte di quelle di test ($\approx 0{,}60$). La tesi attribuiva il divario
all'overfitting e la sconfitta di DINO alla scarsità di dati di pre-addestramento (circa
\num{11000} ritagli contro i milioni di immagini di ImageNet).

\section{Lettura critica}
\label{sec:lettura-critica}

Il confronto tra il testo, il codice originale e i numeri riportati fa emergere diversi
fattori, \emph{alternativi} alle spiegazioni proposte, che possono aver prodotto il risultato
negativo (\cref{tab:lettura-critica}). Il punto per questo lavoro è che le ipotesi di ricerca
\textbf{non sono state confutate in modo conclusivo}.

\begin{table}[t]
  \centering
  \small
  \caption[Osservazioni della tesi del 2022 e fattori alternativi]{Osservazioni della tesi del
    2022, spiegazioni proposte allora e fattori alternativi individuati nel codice e nel
    protocollo.}
  \label{tab:lettura-critica}
  \begin{tabular}{|p{2.6cm}|p{2.6cm}|p{7.2cm}|}
    \hline
    \textbf{Osservazione} & \textbf{Spiegazione proposta} & \textbf{Fattore alternativo} \\
    \hline
    Clutter quasi mai riconosciuto & eterogeneità della classe &
      La perdita escludeva l'indice 0, che corrispondeva proprio al clutter: per quella classe
      il modello riceveva solo gradienti che ne abbassavano il logit
      (\cref{prop:classe-esclusa}). \\
    \hline
    Validation $\gg$ test & overfitting &
      L'overfitting si manifesta come training $\gg$ validation, non validation $\gg$ test. Un
      salto da 0,96 a 0,60 suggerisce una validation non indipendente (ritagli sovrapposti o
      split non disgiunti per area, \cite{kaufman2012leakage}) o metriche calcolate in modo
      diverso (media per batch, \cref{sec:metriche}). \\
    \hline
    DINO $<$ ImageNet & dati di pre-addestramento insufficienti &
      L'implementazione si discostava da DINO: nessuna \emph{projection head}, softmax
      calcolata sull'intero batch appiattito (le distribuzioni dei campioni si mescolano),
      viste locali a 224 invece che a 96 pixel, temperatura del teacher fissa, nessuno schedule
      del \emph{weight decay}, teacher non inizializzato come copia dello student. \\
    \hline
    VQ peggiora & la quantizzazione non aiuta &
      Il codebook quantizzava solo il collo di bottiglia, aggirato dalle connessioni di salto
      (\cref{prop:aggiramento}). Lo stato EMA veniva ricreato a ogni passo. Un encoder
      condiviso per errore tra istanze del modello faceva sì che congelare il teacher
      congelasse anche lo student. \\
    \hline
    ResNet-18 $\approx$ ResNet-34 & basta il modello piccolo &
      Nel codice disponibile non c'è una variante ResNet-34: il codice che ha prodotto le
      tabelle non è con certezza quello conservato. \\
    \hline
  \end{tabular}
\end{table}

A questi fattori si aggiungono alcune scelte di protocollo che impediscono il confronto con la
letteratura.

\begin{itemize}
  \item \textbf{Split non standard}: quattro delle aree rimosse (\tile{2\_13}, \tile{2\_14},
  \tile{6\_13}, \tile{7\_13}) sono tile di \emph{test} nello split più usato, e il testo non
  descrive da dove provengano i \enquote{Test Scores}.
  \item \textbf{Ingressi non normalizzati} con le statistiche di ImageNet, pur partendo da
  pesi ImageNet.
  \item \textbf{Batch di training ordinati} (nessun rimescolamento), quindi spazialmente
  correlati.
  \item \textbf{Una sola esecuzione per configurazione}, senza seed: differenze di uno o due
  punti non sono distinguibili dalla variabilità tra esecuzioni~\cite{bouthillier2021variance}.
\end{itemize}

\section{Implicazioni per questo lavoro}
\label{sec:implicazioni}

La \cref{tab:correzioni} riassume come ciascun problema viene affrontato nella nuova
implementazione e dove la scelta è discussa. Le correzioni non cambiano le domande di ricerca:
le rendono rispondibili.

\begin{table}[t]
  \centering
  \small
  \caption[Problemi del protocollo 2022 e correzioni adottate]{Problemi del protocollo 2022 e
    correzioni adottate in questo lavoro.}
  \label{tab:correzioni}
  \begin{tabular}{|p{4.4cm}|p{6.0cm}|c|}
    \hline
    \textbf{Problema} & \textbf{Correzione} & \textbf{Dove} \\
    \hline
    Clutter escluso dalla perdita & si ignora solo il valore 255 (pixel senza etichetta) & \cref{sec:formulazione} \\
    \hline
    Metriche mediate per batch & matrice di confusione globale sull'intero split & \cref{sec:metriche} \\
    \hline
    Validation non indipendente & split per tile intere; validation su tile escluse dal training & \cref{sec:split} \\
    \hline
    Split non standard & split di riferimento della letteratura; protocollo 2022 mantenuto solo per il confronto & \cref{sec:split} \\
    \hline
    Valutazione su ritagli & finestra scorrevole su tile intere, etichette piene ed \emph{eroded} & \cref{sec:finestra} \\
    \hline
    DINO non fedele & head di proiezione, perdita per campione, multi-crop $2\times224 + 6\times96$, schedule completi & \cref{sec:dino} \\
    \hline
    Codebook fragile & EMA in variabili persistenti, inizializzazione dai dati, riavvio dei codici morti, monitoraggio della perplexity & \cref{sec:vq} \\
    \hline
    Collo di bottiglia aggirato & quantizzazione opzionale anche delle connessioni di salto & \cref{sec:vq-aggiramento} \\
    \hline
    Normalizzazione assente & normalizzazione ImageNet come primo strato del modello & \cref{sec:addestramento} \\
    \hline
    Un'esecuzione, nessun seed & tre seed per configurazione, media $\pm$ deviazione standard & \cref{sec:statistica} \\
    \hline
  \end{tabular}
\end{table}
